\section{Method families and their medical evidence}
\label{sec:methods}

Following \S\ref{sec:taxonomy}, each subsection defines a family, reviews representative
studies, tabulates the medical evidence and summarises its claims and limits; studies with claims
in two families appear in both. Here ``medical'' means analysing a medical foundation model
(\S\ref{sec:models}); cited general-domain studies supply methods or mechanisms used by medical work. Table~\ref{tab:families} summarises what each family
shows and the tests that strengthen it, and Fig.~\ref{fig:families} draws the ladder of evidence from
observation to concept-specific causal contribution that the formal view below defines. Core studies whose quality note records a defect that undermines the representation-level result (\S\ref{sec:findings}; the judgement is recorded per record in the supplement and counted per finding in supplementary Tables~S9--S11) are marked $^{\dagger}$ where cited singly in this section (a group citation is not marked), so that their results are read with that qualification.


\textbf{Formal view.} In the notation of \S\ref{sec:prelim}, each family can be read as characterising a
functional of the representation $z_\lambda(x)$ at a locus $\lambda$, differing in the functional
examined and in whether the forward pass remains intact. For a multimodal model the state also depends on the prompt $t$, the token position $p$ and, for generated
end-points, the decoding policy, so $z_\lambda(x)$ abbreviates $z_\lambda(x,t,p)$ and $B(x)$ abbreviates
$B(x,t)$ under a fixed decoding rule. A claim about a decoder or connector state is therefore conditional on
the prompt and position at which it was measured, which is why the coding records locus and readout
separately. Structural analysis reports a
statistic $\Phi(\{z_\lambda(x_i)\}_{i=1}^{N})$ of the embedding distribution, unconditionally
or stratified by metadata such as scanner or site, $\Phi(\{z_\lambda(x_i)\}_{i:\,m(x_i)=m})$;
neither decodes a concept. Decoding fits, within a decoder class $\mathcal{G}$ and on a training set
$\mathcal{T}$, the map
\begin{equation}
\hat g=\arg\min_{g\in\mathcal G}\ \tfrac{1}{|\mathcal T|}\textstyle\sum_{i\in\mathcal T}\ell\big(g(z_\lambda(x_i)),\,c(x_i)\big)
\label{eq:probe}
\end{equation}
and reports the performance of $\hat g$ on examples held out of $\mathcal T$, which depends on
$\mathcal G$ as much as on $z_\lambda$. Discovery fits a dictionary $D$ and sparse codes $a(x)$ with $z_\lambda(x)\approx D\,a(x)$,
then names its columns. Tracing reports functionals of the unchanged forward pass indexed by
layer, position or component, such as layer-wise decoding of $h^{(\ell)}$ or attention mass on
visual tokens. Intervention replaces $z_\lambda$ by $\tilde z_\lambda$ and reports the
behavioural effect
\begin{equation}
\Delta B=B\big(x;\,z_\lambda\!\leftarrow\!\tilde z_\lambda\big)-B(x),
\label{eq:effect}
\end{equation}
with $\overline{\Delta B}=\mathbb{E}_x[\Delta B(x)]$ its expectation over inputs (and, for generated text, over
the decoding policy), the replacement $\tilde z_\lambda$ being in general itself a function of $x$, $t$ and
$p$; the reviewed studies report sample averages of it. Equation~(\ref{eq:effect})
describes within-run interventions; post-hoc transformations of stored embeddings, parameter editing
and unlearning have their own estimands and subtype-specific end-points (\S\ref{sec:intervention}).
For a steering intervention $\tilde z_\lambda=z_\lambda+\alpha v_c$ along a direction built for concept
$c$, the coding records whether the study reports any matched control against which $\Delta B(\alpha v_c)$
is compared; \S\ref{sec:future} proposes a diagnostic that operationalises the comparison, as a reporting
proposal rather than a standard the reviewed literature was held to. Exceeding matched controls establishes
selective sensitivity to the direction, not semantic validity, which depends on how $v_c$ was obtained.
Erasure carries part of its control in its construction, an affine map removing the linear predictability
of $c$ under a stated metric~\cite{belrose2023leace}, but attributing a behavioural change to it still
requires a sham or rank-matched transformation, and checks that $c$ is not recoverable non-linearly, that
correlated concepts survive and that off-target performance is retained. The paired probe-and-intervention design that the medical corpus lacks (\S\ref{sec:intervention}, \S\ref{sec:future}) follows directly: $v_c$ is the normal of a binary linear probe fitted by~(\ref{eq:probe}) or a dictionary column decoding $c$, and~(\ref{eq:effect}) is reported against matched control directions at the same intervention locus and behavioural end-point.

Eq.~(\ref{eq:probe}) is the canonical decoding case; prompt-similarity scoring recovers $c$ through a fixed similarity structure rather than a fitted $\hat g$, and labelled representational similarity shows only that representation distances covary with concept-labelled distances, without mapping a held-out example to a concept value; both are grouped with decoding by the claim they license, the second as the weakest member (\S\ref{sec:taxonomy}). \textbf{Three causal claims.} Three claims of increasing strength follow from~(\ref{eq:effect}).
\emph{Causal sensitivity} means that the output changes when $z_\lambda$ is changed; any
intervention whose effect exceeds its uncertainty over inputs establishes it. \emph{Direction-selective sensitivity}
means that the effect of $v_c$ exceeds that of matched controls (the screening diagnostic proposed in \S\ref{sec:future}). \emph{Concept-specific causal contribution} means that the effect is attributable to the model's representation of $c$ rather than to a merely correlated direction. Necessary but not sufficient conditions for it, beyond matched controls, are validation of $v_c$ through
agreement among several independently constructed directions for the same concept, a
correlated-concept hard-negative control, and an intervention that stays within the empirical activation distribution (the layer's range of activation norms and its principal subspace --- a proximity check, not a test of manifold membership).
Edits far from that distribution can change behaviour for reasons unrelated to $c$; together these checks falsify the simplest alternatives and establish stronger direction-selective sensitivity, consistent with rather than proof of a concept-specific contribution. Identifying it takes a mediation design with a defined exposure, mediator and outcome and the sequential-ignorability assumptions that identify a natural indirect effect~\cite{pearl2001direct,imai2010mediation}, a causal-concept effect estimated under an intervention on the concept itself~\cite{goyal2019cace}, or a causal-abstraction mapping between concept and representation~\cite{geiger2021causal,geiger2023causal}; no medical study in this corpus reaches any of the three. The gap between structure that a representation merely exhibits, structure that an intervention can move, and factors that are causally meaningful is the one Sch\"olkopf \emph{et al.}~\cite{scholkopf2021causal} draw for representation learning generally; the three inferential levels of this review are its measurement-side counterpart. Table~\ref{tab:families} (\S\ref{sec:comparison}) restates this view family by family.

\begin{figure*}[!tbp]
\centering
\includegraphics[width=0.95\linewidth]{fig6.png}
\caption{From what is encoded to what is used. Observational evidence --- structure, decodability,
discovered features, localisation --- characterises $z_\lambda(x)$ and shows availability, not use. Use
requires an intervention $z_\lambda\to\tilde z_\lambda$ with a pre-specified behavioural end-point
$\Delta B$, and the interventional evidence then climbs three rungs of increasing inferential strength:
\emph{causal sensitivity} ($\Delta B\neq 0$ on held-out inputs; the edited direction may still be
arbitrary), \emph{direction-selective sensitivity} (the effect of $\alpha v_c$ exceeds equal-norm random
and unrelated-concept controls; still not proof of concept-specific causality) and \emph{concept-specific
causal contribution}, for which semantic agreement, hard negatives and an intervention within the
empirical activation distribution are necessary but not sufficient: identifying the contribution takes a
mediation, causal-concept or causal-abstraction design with stated assumptions (\emph{Formal view}), and
no reviewed medical study meets the criterion. The required evidence grows from a behavioural end-point,
to matched controls, to semantic and distributional validation and, finally, an identification design.}
\label{fig:families}
\end{figure*}

\subsection{Structural and geometric analysis}
\label{sec:geometry}

Structural analyses characterise representation spaces through image--text arrangement,
distributional spread, cross-model similarity, organisation by scanner or site and change under
adaptation. Fitting no decoder, they are the cheapest analyses. Eligibility here requires a demonstrated
link to decodability, retrieval, transfer or shortcut behaviour~\cite{geirhos2020shortcut}, which makes this family a
\emph{functionally grounded} selection rather than a map of medical foundation-model geometry: a study that
measures anisotropy or alignment and reports no downstream consequence does not enter, so purely descriptive
and negative geometry results are under-represented by construction. The restriction is not applied to the
other families, so geometry counts and positive-direction rates are not comparable with theirs, and the
statements it supports (F3c, F5, F5b, F11) should be read as statements about functionally grounded geometry
work.

\subsubsection{Modality gap and cross-modal alignment}

\textbf{The modality gap.} Contrastive dual encoders separate image and text embeddings; the gap
between modality centroids is $\Delta=\|\bar z_v-\bar z_t\|_2$, with $\bar z_v=\frac{1}{N}\sum_i z_v(x_i)/\|z_v(x_i)\|_2$ and $\bar z_t$ defined the same way: each embedding is unit-normalised \emph{before} averaging, so $\Delta$ is a distance between mean directions and $\|\bar z_v\|_2\le 1$; averaging first and normalising afterwards gives a different quantity, and the studies that report $\Delta$ are not always explicit about which they compute. General-domain
accounts trace it to initialisation cones, information imbalance, sparse cross-modal atoms or a
functional role~\cite{liang2022mind,schrodi2025two,dhimoila2026crossmodal,yan2025beyond}. Restrepo
\emph{et al.}~\cite{restrepo2026cone} find tighter cones on medical than on natural data for all four backbones and a
smaller cross-modal separation in BiomedCLIP and MedSigLIP than in CLIP and SigLIP, though the gap persists, and a frozen-encoder
shift gives task-dependent gains below 0.02 in the area under the ROC curve (AUROC) without a random-direction control. Gap-shrinking losses and gaze guidance improve retrieval, captioning and
alignment~\cite{grassucci2025closing,lahmar2026bridging,kumar2024improving}, on no shared protocol.

\textbf{Alignment quality and cross-modal consistency.} Contrastive objectives decompose into
alignment and uniformity~\cite{wang2020understanding}: paired distance $\mathbb{E}_{(x,t)}\|z_v(x)-z_t(t)\|_2^{2}$ and spherical coverage
$\log\mathbb{E}_{x,x'}\exp(-2\|z(x)-z(x')\|_2^{2})$, though intra- and cross-modal similarities disagree in CLIP~\cite{mistretta2025cross}. Global
alignment survives where local alignment does not: in GLoRIA, global image--sentence similarity separates the
paired sentence from a random one at 70\% but region-level similarity only at 55\%, and at 43\% on abnormal
pairs, below chance~\cite{mcinerney2022that}; counterfactual edits of negation and laterality barely move
BiomedCLIP and PubMedCLIP similarity~\cite{balogh2026regularization}$^{\dagger}$; and a projection head can
collapse tumour and normal prototypes to 98\% similarity~\cite{arora2026semantic}. Because distinct patients
share findings, off-diagonal negatives are often semantically positive, and the objectives addressing
this~\cite{wang2022medclip,park2025off,zhang2025fane,ztrk2025meta} are evaluated downstream rather than
geometrically.

\subsubsection{Anisotropy, effective rank and dimensional collapse}

The cone effect, anisotropy, low effective rank and dimensional collapse all describe concentration
in a small subspace, first found in language embeddings~\cite{ethayarajh2019how}; rapid spectral decay
indicates low effective rank, $\mathrm{erank}=\exp(-\sum_k p_k\log p_k)$ with
$p_k=\sigma_k/\sum_j\sigma_j$ the singular values of the centred embedding matrix normalised to sum to
one, so $p_k$ is a share of singular value and not of variance (using $\sigma_k^2$, equivalently the
covariance eigenvalues, gives a smaller number, and the studies compared here are not always explicit about
which they report); effective rank predicts downstream performance; how uniformly a space is used can be measured
directly~\cite{rudman2022isoscore}, and contrastive training can collapse
dimensions~\cite{roy2007effective,garrido2022rankme,jing2021understanding,hassanpour2024overcoming}. Across
eight encoders, aggressive post-hoc whitening (\S\ref{sec:intervention}) lowered downstream AUC in
34 of 36 model--dataset pairs, though the pooled interval crossed zero, whereas gentler spectral
standardisation gained about 0.005 AUROC (from the abstract)~\cite{garciahidalgo2026embedding}. Since the
near-isotropic encoder degrades with the rest, that experiment does not establish the anisotropy itself as
the task-relevant structure, and F5 records the point as a single-study hypothesis. Concentration
neither places concept directions in the dominant subspace nor precludes decodability along
low-variance directions, so similarity scores need standardisation and, with
superposition~\cite{elhage2022toy}, overcomplete dictionaries (\S\ref{sec:discovery}).

\subsubsection{Cross-model similarity}

Cross-model similarity is measured by linear centred kernel alignment
(CKA)~\cite{kornblith2019similarity}, representational similarity analysis
(RSA)~\cite{kriegeskorte2008representational}, canonical-correlation
measures~\cite{raghu2017svcca,morcos2018insights} or Procrustes distances, whose sensitivities differ and are
surveyed and benchmarked elsewhere~\cite{klabunde2023similarity,ding2021grounding}: for centred activations $X$
and $Y$, $\mathrm{CKA}(X,Y)=\|Y^{\top}X\|_F^{2}/(\|X^{\top}X\|_F\,\|Y^{\top}Y\|_F)$, while RSA correlates pairwise-dissimilarity matrices. Under the
platonic-representation hypothesis~\cite{huh2024platonic}, visual tokens align with text through the
stack~\cite{neo2024towards}, whereas pathology encoders track slide of origin over
disease~\cite{mishra2025comparing}.

The largest medical study examines 18 image and 7 text encoders across five modalities in an
objective-stratified training matrix (one seed per cell, native preprocessing)~\cite{tayebiarasteh2026self}:
convergence is modest and groups by objective rather than scale, with matched self-supervised encoders
aligning most (40.4\% on chest radiography) against label-supervised (21.1\%) and image--text (3.3\%) ones.
No image encoder aligns with a clinical text encoder, although linear classifiers transfer at about 85\%. Across 22 encoders, 20 datasets
and 605,443 images, chance-calibrated nearest-neighbour discordance puts blind sets above a
hypergeometric null for every dermatology, fundus, histopathology and mammography pair and 655 of
660 chest pairs, against about half of natural-image pairs~\cite{arasteh2026candor}; nearest-neighbour
statistics in high-dimensional spaces are themselves distorted by hubness~\cite{radovanovic2010hubs}, which
no reviewed study corrects for. Elsewhere between-scanner reliability of brain encoders is poor, CKA falls under MRI artefacts while
effective rank holds, nearest-neighbour overlap among volumetric encoders tracks model quality, and an
at-most-affine map stitches frozen chest-radiograph encoders into a zero-shot
classifier~\cite{navarrogonzalez2026crossscanner,mielcarz2026do,aerts2025foundation,folco2025semantic}
(\S\ref{sec:intervention}). The corpus's cross-encoder
comparisons group by objective~\cite{tayebiarasteh2026self} and by modality~\cite{arasteh2026candor}, not by
domain (\S\ref{sec:findings}). None of this is surprising if pipelines are underspecified: equal downstream
performance does not identify equivalent representations, and the choices left free by the training
objective can move them~\cite{damour2022underspecification}; identifiability up to a linear map holds only under
conditions on the training objective that these pipelines do not check~\cite{roeder2020linear}.

\subsubsection{Site, scanner and acquisition structure}

Acquisition confounding predates foundation models~\cite{zech2018variable}, and medical embeddings
are organised the same way
(\S\ref{sec:decoding}~\cite{dejong2025current,komen2024do,zheng2024demographic}). The signatures
behave as geometric distortions: one affine correction per model harmonises embeddings across
pathology scanners and MRI field strengths~\cite{donle2026featmap}, scanner variability shifts calibration
while preserving discrimination~\cite{thiringer2026scanner}, and a contrastive projection on frozen
embeddings reduces scanner clustering~\cite{carloni2025pathology} (\S\ref{sec:intervention}).

Pathology provenance is recoverable at high accuracy: source hospital from ten encoders at AUC 0.82--0.99,
site above 0.95 after Macenko normalisation, scanner from slide embeddings above 0.99 out of fold, and
14--55\% of a patch embedding's neighbours reflecting centre rather than
biology~\cite{zhang2025hospital,murchan2024deep,kmen2026robust,jung2026foundation,shafique2026investigating,saueressig2025histology,gustafsson2024evaluating,esbri2025benchmarking}.
The consequence is not only detectability: the scanner signature covaries with case mix and inflates an
external Ki-67 correlation~\cite{jung2026foundation}, stain normalisation, lossy DICOM conversion and
sequential sectioning all move embeddings and
predictions~\cite{wolflein2023benchmarking,mayer2025lossy,balan2025reproducibility}, and in CT and PET
embedding concordance holds across reconstruction settings while nodule-outcome ranking moves with
them~\cite{park2025evaluating,farquhar2026pretraining}.

Outside pathology, a t-SNE of frozen Mammo-CLIP features visually suggests clustering by dataset across ten mammography cohorts, with no separation statistic reported, and
the encoder most organised by cohort probes best yet loses most under external transfer, ending level with a general
encoder~\cite{germani2025bias}. Six further studies tie embedding structure to subgroup and imaging attributes, judged image quality,
tissue-type clustering, augmentation, biological order and institutional
artefacts~\cite{alloula2025representation,ouaari2026assessing,fournier2026prognostic,elphick2024are,vig2026do,zhang2026auditing}.
Image--report correspondence in a chest-radiograph joint space, at 15--50 times chance, is a
re-linkage risk~\cite{tayebiarasteh2026crossmodal} (\S\ref{sec:tracing}), motivating the confounder
controls of \S\ref{sec:applications}.

\subsubsection{Representation change under adaptation}

Similarity, checkpoint probes and cross-model dictionaries measure change under fine-tuning,
low-rank adaptation and continued pre-training. General-domain work finds that fine-tuning most changes top layers and can distort features, that
low-rank updates introduce intruder dimensions, and that visual instruction tuning moves visual features
into intermediate language-model
layers~\cite{merchant2020what,zhou2021closer,kumar2022fine,lee2022surgical,raghu2019rapid,biderman2024lora,shuttleworth2024lora,palacios2026visual,verma2024cross}.
Medical few-shot and meta-learning are assessed downstream~\cite{bo2025famnet,wang2025task}, and the
representation-level studies measure cross-hospital drift, forgetting during continued pre-training and
adapter placement from transferred
dictionaries~\cite{theofilou2025stable,zoellin2024block,sadanandan2026mechanistically}. Whether concept
directions survive transfer is unmeasured, although crosscoder model diffing~\cite{minder2025overcoming} is
an instrument (\S\ref{sec:discovery}).


\subsubsection{Summary and discussion}

Medical encoders converge modestly, grouping by objective rather than scale~\cite{tayebiarasteh2026self},
were alignable in one study by an affine map to within about a dozen AUROC points of the source
encoder~\cite{folco2025semantic}, and place acquisition variables in prominent directions.
\Vgeoany{} of the \Vgeon{} geometry records compare the statistic with a comparison model or a null
(supplement S5d); whether geometry-modifying methods repair what they target is unmeasured,
because they are assessed only downstream.

\subsection{Supervised concept mapping and decodability}
\label{sec:decoding}

Decoding maps representations to external concepts through probes, concept activation vectors,
labelled-set similarity or zero-shot text scoring. The oldest family, it measures decodability, a joint
property of representation and decoder~\cite{alain2016understanding,belinkov2021probing,hewitt2019designing,elazar2020amnesic}.

\subsubsection{Probing for clinical findings and anatomy}

Seventeen frozen-feature benchmarks meet the criteria of \S\ref{sec:method}, \Nbenchdec{} of
them with decoding-family records and two with structural records only. They rank
frozen pathology, chest-radiograph, CT, retinal and volumetric
encoders~\cite{neidlinger2025benchmark,marza2025thunder,li2025featurequality,li2025embeddings,chevli2026frozenct,zhou2025generalist,aerts2025foundation}.

Across ten mammography cohorts, frozen Mammo-CLIP probes beat GLoRIA, CLIP and DINOv2 by 9.7 to 15.3\% relative F1
averaged over the diagnosis and density tasks, the diagnosis F1 being about 0.65 against 0.56--0.58, yet a probe fitted
on a single cohort falls from 0.73 to 0.32 on the other nine~\cite{germani2025bias}. On 492,724 chest radiographs, pooled embeddings from five frozen
transformers are at chance for the controlled small-scale perturbations (AUC 0.50--0.52) while patch tokens
from the same forward pass recover them; on the 3,543-image bounding-box cohort the same contrast holds, with
patch-local AUC at least 0.899 in every model--class cell, but the patches are selected by the ground-truth
lesion box, so the recovery is an oracle readout rather than a deployable one~\cite{muthyala2026frozen}. On a
pre-registered PET benchmark of 2,592 patients, patch classification saturates for every frozen
encoder and for a randomly initialised control~\cite{farquhar2026pretraining}.

The same readout reaches phenotype codes, training configurations, knee grades, cell types, omics tasks,
spatial embedding tensors, placental lesions, prostate grade, spatial-domain labels, kidney morphometrics and
radiological retrieval, in benchmarks of five to nineteen
encoders~\cite{lin2026transparent,liang2026cheap,demir2025downstream,cui2026translating,hu2026probing,jamzad2026p3ca,sadman2025interpreting,svatko2026deep,jung2026foundation,steiner2026when,vadori2025mind,bakhta2026frozen,beyeler2024comparing,peng2025benchmarking,gustafsson2024evaluating,kasireddy2025explainable,denner2024leveraging,howard2024generative,ramchandani2026benchmarking,zhao2026benchmarking}
Two of them carry a critical-defect flag, for patch-level leakage and a
circular cluster-selection metric, and the instructive case is a failure: an ejection-fraction probe with
$R^2$ 0.63 in the development cohort correlates negatively with the label on an external
one~\cite{akanova2026interpreting}$^{\dagger}$.

Decodability can also be localised, to intermediate layers for chest-radiograph shortcuts, to a dictionary
neuron for tumours, and to roughly ten channels whose ablation removes a finding
altogether~\cite{pedersen2026look,tamura2025patch,nooralahzadeh2026sparse,goetze2025prediction,basu2026interpretability,nooralahzadeh2026arbitration}.

\subsubsection{Probing for demographic and acquisition attributes}

Before foundation models, chest-radiograph classifiers underdiagnosed underserved groups, failed on hidden
strata and acquisition cues, predicted race, sex and age without anatomical correlates, and showed
misdiagnosis disparities in
pathology~\cite{seyyedkalantari2021underdiagnosis,oakdenrayner2019hidden,degrave2020ai,banerjee2021ai,adleberg2022predicting,korot2021predicting,coyner2021color,burns2023ability,vaidya2024demographic,jones2024causal};
asking whether an attribute is encoded or used made this a representation-level
question~\cite{glocker2023algorithmic,glocker2022risk}.

Frozen CT embeddings trained without demographic labels decode sex at AUROC 0.998, age at RMSE 3.8 years and
race at 0.878~\cite{zheng2024demographic}, and the pattern recurs in chest radiography, pathology, brain MRI,
retina and mammography, where embeddings also give away site, camera or dataset even after stain
normalisation~\cite{sourget2025fairness,dejong2025current,komen2024do,rahbar2026frozen,li2026acquisition,hajishafiezahramini2026dataset,encin2026comparative,lin2025beyond,kmen2026robust}.
The range is wide, and the low end is instructive: eight chest-CT encoders decode race at only 0.72--0.74
AUROC and read age and sex no better than a randomly initialised
encoder~\cite{tariq2024self}$^{\dagger}$, while of thirteen shortcut detectors on CheXpert only the non-probe
ones flag race~\cite{cajas2026shortkit}.

On 230,697 radiographs, sex is encoded at AUROC 0.942, yet linear finding probes refitted after residualising
it out lose less than 0.001 AUROC, so the sex component is redundant for an optimally refitted linear readout;
whether the model's own head used it was not tested~\cite{kim2026encoding}$^{\dagger}$. Orthogonalisation recovers these attributes before projection and chance after, cropping reduces them, and
fairness benchmarks span 17 datasets and 20
models~\cite{weber2023posthoc,sutariya2026impact,jin2024fairmedfm,kauffmann2024explainable}.

\subsubsection{Concept activation vectors and directional sensitivity}

Concept activation vectors (CAVs) fit a direction from labelled positive and negative examples, then
score classes by logit directional derivatives~\cite{kim2018interpretability}. With $v_c$ the unit
normal of a classifier separating $\{z_\lambda(x):x\in P_c\}$ from negatives and $\eta_k$ the model's logit
for class $k$, sensitivity is $S_{c,k}(x)=\nabla_{z_\lambda}\eta_k(x)\cdot v_c$ and the TCAV score is
the fraction of class-$k$ inputs with $S_{c,k}(x)>0$, tested against random directions. Applications cover histopathology, dermoscopy, cardiac MRI and diabetic-retinopathy
grading~\cite{graziani2018regression,graziani2020concept,lucieri2020interpretability,clough2019global,stors2024looking},
and iterated HPV directions predict HPV status in a pathology
encoder~\cite{hieromnimon2025analysis}$^{\dagger}$. Text-derived vectors remove labelled exemplars and support
adversarial debiasing~\cite{nicolson2024textcavs,correa2023efficient}, but the scores stay observational, and
they are inconsistent across layers, entangled and spatially dependent~\cite{nicolson2024explaining}.

\subsubsection{Representational similarity with concept-labelled sets}

Representational similarity analysis compares representation and stimulus distances without a
decoder~\cite{kriegeskorte2008representational}. Across six pathology encoders similarity follows
slide identity more than disease, which stain normalisation only partly
reduces~\cite{mishra2025comparing}. Diffusion pseudotime over five pathology encoders recovers the
progression order of four cancer types against a label-shuffled null, order fidelity predicting
few-shot accuracy~\cite{vig2026do}. Chance-calibrated discordance across 22 encoders, anchored by a random encoder, shows one encoder reading
pneumothorax at 84.5 AUROC while placing 18\% of positives nearer an opposite-label film from the same
hospital~\cite{arasteh2026candor}.

\subsubsection{Text-anchored concept scoring}

Prompt similarity in joint spaces gives zero-shot concept scores. MONET audits data and models with
dermatologist-verified concepts~\cite{kim2024transparent}, which pre-training on 23 million
image--text--concept triplets extends~\cite{nie2025explainable}. The same zero-shot scoring serves as a general-purpose concept detector, for dermoscopic criteria, retinal
clusters, concept-to-region grounding and counterfactual
generation~\cite{patricio2023conceptbased,mensah2026interpretable,luo2025xbench,sadashivaiah2024explaining}, with one
instructive limit: it detects injected biases but not implicit ones~\cite{bertolini2025strengths}. Because prompts alter directions,
expert-label validation is required (\S\ref{sec:eval-what}).

\subsubsection{Validity controls}

Expressive decoders require controls: a control task refits $\mathcal G$ in~(\ref{eq:probe}) on permuted or
type-matched labels and reports task-minus-control accuracy~\cite{hewitt2019designing}, description length
controls capacity~\cite{voita2020information}, amnesic probing and erasure test whether removing a property
changes behaviour~\cite{elazar2020amnesic,ravfogel2020null,belrose2023leace}, and information-theoretic and
geometric accounts bound cross-representation comparison and identify concept
subspaces~\cite{pimentel2020informationtheoretic,guerner2023geometric}. Medical use is rare. Among the \Ndecni{} core decoding studies outside intrinsic
designs, \Nprobectrl{} report a shuffled- or permuted-label null or a chance
calibration~\cite{liang2026cheap,rahbar2026frozen,vig2026do,bertolini2025strengths,arasteh2026candor,steiner2026when,li2026acquisition,cajas2026shortkit},
none a randomized control task in the sense of
Hewitt and Liang (\S\ref{sec:findings}); \Nproberand{} report a randomly initialised
encoder~\cite{rahbar2026frozen,arasteh2026candor,svatko2026deep,farquhar2026pretraining} and \cite{tariq2024self}$^{\dagger}$, two of them also among the first (\S\ref{sec:eval-what} gives the full partition).
Twenty-eight compare medical and general models on the same data. Only four probe successive depths in order to
localise where a concept becomes decodable~\cite{pedersen2026look,rahbar2026frozen,vig2026do,vadori2025mind};
two more read several depths without localising~\cite{farquhar2026pretraining,demir2025downstream}, one varies
pooling rather than layer~\cite{muthyala2026frozen}, and the rest probe the final embedding. A frozen-probe benchmark of physiological-signal models sets the standard~\cite{zare2026eeg}, and the
interventions of \S\ref{sec:intervention} supply the behavioural end-points decoding lacks.


\subsubsection{Summary and discussion}

Whether a recoverable attribute influences the output is usually untested: only
residualisation~\cite{kim2026encoding}$^{\dagger}$ and channel ablation~\cite{nooralahzadeh2026sparse}
separate availability from the dependence of a specified downstream readout, and neither tests the model's
own fixed head. Decodability of a named concept across successive loci of one model, under a control, remains unmeasured;
the one successive-locus probe reads class labels without one (F8b, \S\ref{sec:tracing}).

\subsection{Feature decomposition and discovery}
\label{sec:discovery}

Discovery decomposes representations into neurons, factorised directions or sparse features without a
prior vocabulary, then names them from activating inputs; naming validity and stability decide
what it establishes.

\subsubsection{Neuron-level analysis and network dissection}

Network dissection labels each unit with the concept whose mask maximises intersection over union with
its activation map~\cite{bau2017network}, CLIP-Dissect substituting a vision--language
model~\cite{oikarinen2022clipdissect}; the multimodal neurons of CLIP, responding to a concept whether
it is drawn, photographed or written, are the precursor of unit-level multimodal discovery~\cite{goh2021multimodal}. Mid-layer atlases of a pathology encoder show coherent tissue
concepts but dispersed subclasses, pathologist agreement falling from $\kappa$ 0.82 for cancer groups to
0.11 for subclasses~\cite{gustav2026class}; other unit-level work labels mammography neurons, rebuilds
a diagnostic model, probes retinal networks and reads hallucination off language-model
units~\cite{salahuddin2025mammoclip,gong2025concepts,yousefzadeh2023lava,vankadaru2026readable}.
Superposition and non-axis-aligned concepts limit unit
analysis~\cite{elhage2022toy,park2023linear}, a concept is in general carried by a combination of
filters rather than by one~\cite{fong2018net2vec}, and unsupervised factors are not identifiable without
inductive bias or supervision~\cite{locatello2019challenging}, so medical work benchmarks dictionaries (Table~S13).

\subsubsection{Factorisation directions}

Principal components, independent components and non-negative factors give few directions, which in
medical embeddings reflect acquisition and site before pathology (\S\ref{sec:geometry}): among inverted
principal components of pathology-encoder gradients, one stain-colour direction carries 54--69\% of the
tissue-source-site gradient against at most 20\% for any grade component~\cite{howard2024generative}.
Other factorisations recover plausible morphology: concept-bank segmentation beats a
ResNet-50~\cite{gildenblat2024segmentation}$^{\dagger}$ and three pathologist-named HPV-positive
subtypes emerge~\cite{hieromnimon2025analysis}$^{\dagger}$. Carcinoma morphology, retinal characteristics,
breast-cancer archetypes and an infectious-disease proposal are reported
alongside~\cite{qin2026clearhpv,varshney2024generating,beyeler2024comparing,fournier2026prognostic,sizikova2022automatic}.
Activation clustering~\cite{ghorbani2019automating,fel2023a} and sparse decomposition onto a fixed
vocabulary~\cite{bhalla2024interpreting} have no medical instance; dictionary learning grew from
monosemanticity work~\cite{bricken2023monosemanticity,templeton2026scaling}.

\subsubsection{Sparse autoencoders and dictionary learning}

A sparse autoencoder (SAE) reconstructs activations through an overcomplete, sparsely active layer of
candidate monosemantic features~\cite{cunningham2023sparse}. With codes
$a(x)=\sigma(W_e z_\lambda(x)+b_e)$ and reconstruction $\hat z=W_d\,a(x)+b_d$, the dictionary
minimises
\begin{equation}
\begin{split}
&\mathbb{E}_x\big[\|z_\lambda(x)-\hat z(x)\|_2^{2}+\beta\,\|a(x)\|_1\big]\\
&\quad\text{s.t.}\ \|(W_d)_{:,j}\|_2=1\ \forall j,
\end{split}
\label{eq:sae}
\end{equation}
where $(W_d)_{:,j}$ is the $j$th column of the decoder, the direction of feature $j$.
Equation~(\ref{eq:sae}) is the canonical $\ell_1$ objective; the reviewed medical dictionaries replace the
penalty by a fixed active-code count (top-$k$, batch-top-$k$), a gated or jump activation or a matched-pair
crosscoder loss, changing what is held fixed rather than the estimand. Unit norm removes scale degeneracy,
and each $W_d$ column is a feature named from its highest-code inputs, a naming route that need not identify what the feature does, since in vision the maximally activating patches of a dictionary feature can fail a causal test of the same feature~\cite{han2025causal}. Scaling~\cite{gao2024scaling}, better
activations~\cite{rajamanoharan2024jumping} and open dictionaries~\cite{lieberum2024gemma} made SAEs
portable; medical adoption has two waves. Dictionaries that decompose an image--text space jointly, rather than one tower at a time, are a general-domain development the medical corpus has not taken up: group-sparse autoencoders over paired embeddings~\cite{kaushik2026decomposing} and canonical correlation between sparse concept sets~\cite{gordon2026scocca} both address the cross-modal case that the modality gap of \S\ref{sec:geometry} makes hard, and no reviewed medical study uses either.

The first wave trains on \emph{frozen encoders}. Pathologists curate 41 sarcoma concepts on CONCH
embeddings, where the full dictionary reaches 73.7\% balanced accuracy against 77.4\% for the raw
aggregator and 67.0\% for the curated concepts~\cite{bisson2026unsupervised}$^{\dagger}$; restricting the
classifier to the 41 curated features reduced accuracy, but representation, feature count and curation
change together and leakage is flagged, so the cost of naming itself is not isolated. Others take names
from report-distilled descriptions, tissue structures, pathologists or neuropathologists, radiologist
regions or an automated
labeller~\cite{abdulaal2024an,le2024learning,kim2026dissecting,zhao2026compositional,redlich2026explainable,nakka2025mammosae,renzulli2025medsae},
and they compress: ten features recover most downstream CT and MRI performance~\cite{wesp2026sparse} and
a geometry-guided dictionary predicts Alzheimer's conversion~\cite{nerrise2026geosae}. Further variants
appear in Table~S13~\cite{tang2025cxrlanic,sun2025prototypebased,dasdelen2025cytosae,yan2026a,huang2026classguided,chhetri2026when}.
The second wave trains on \emph{residual streams of medical multimodal LLMs}, recovering device,
pathology and longitudinal-change features from a report generator~\cite{bouzid2025insights}, steering
directions on three radiology models (\S\ref{sec:intervention})~\cite{nooralahzadeh2026universal} and a
transferable general suite~\cite{sadanandan2026mechanistically}; domain-specific dictionaries explain more
clinical-activation variance~\cite{oneill2025resurrecting}. No eligible medical study trains a dictionary
across models, across modalities or on model differences. The templates exist in general
work~\cite{thasarathan2025universal,nasirisarvi2025sparc,minder2025overcoming,plantinga2025from}, in
dictionaries over medically fine-tuned language models, task-specific segmentation models and
physiological-signal
models~\cite{tan2026efficient,ahmed2026medseglens,lehnschiler2026mechanistic,duan2026cadence}, and in the
general-vision dictionaries that steer or segment through the visual
stream~\cite{joseph2025steering,pach2025sparse,stevens2025interpretable}, the designs \S\ref{sec:future}
points to.

Across the \Nsae{} core dictionary studies (supplementary Table~S13 lists every
dictionary study in the corpus), every axis that fixes a feature varies:
token, layer, expansion and sparsity. Features are therefore not commensurable, and naming is as
uneven.
Validation of feature meaning uses clinician readers in \Nsaeclin{} studies and structured labels in
\Nsaelabel{}. Only the \Nsaeval{} with clinician or label validation enter the breadth-labelled statement
F7 (\S\ref{sec:findings}); only \Nsaestab{} report reproducibility, none matches features across seeds
and only \Nsaerecon{} quantify reconstruction, so headline claims of ten-feature
compressibility~\cite{wesp2026sparse} and five thousand monosemantic
features~\cite{tang2025cxrlanic} are neither validated nor reproduced.

\subsubsection{Naming, validation and stability}

Automated pipelines name features from activating inputs and score the
predictions~\cite{paulo2024automatically}; the medical instances grade names against reports, labels or
regions, or discover features and then name
them~\cite{bouzid2025insights,abdulaal2024an,haque2026medconcept,nakka2025mammosae,rao2024discover}, none
testing cross-labeller naming. Stability is worse: absorption and splitting persist across widths and
sparsities~\cite{chanin2024a}, seeds share about 30\% of features~\cite{paulo2025sparse}, derived
concepts are adversarially fragile~\cite{li2025evaluating}, and metrics and benchmarks are too unstable to
separate architectures~\cite{chanin2026are,chanin2026synthsaebench}. Few features reproduce in an
event-sequence anchor, where dictionary probes do not beat dense ones~\cite{sainsbury2026sparse}, and
capacity-matched suites are proposed for signal models~\cite{duan2026ecginterpbench}. Named features are
hypotheses, not stable identifiers, though subspace reconstruction survives.


\subsubsection{Summary and discussion}

Names assigned under weak supervision pass their uncertainty to every downstream claim, and encoder
dictionaries largely recover what probing already sees. Two-seed training, cross-cohort feature matching and
any decomposition at the connector output remain open.

\subsection{Representation attribution, localisation and tracing}
\label{sec:tracing}

Tracing observes forward passes without changing computation, localising concepts to layers,
positions, heads or training examples, each a hypothesis for later patching or steering; the medical
studies ask where multimodal LLMs use images.

\subsubsection{Residual-stream decomposition and lenses}

The logit lens reads an intermediate state through the unembedding matrix,
$\mathrm{softmax}(W_U\,\mathrm{LN}(h^{(\ell)}_p))$; the tuned lens learns a layer-specific affine
map~\cite{belrose2023eliciting}. In vision--language models visual tokens become text-decodable and
middle layers enrich then refine, while visual information declines across layers and generation
steps~\cite{neo2024towards,jiang2024devils,li2025the}; objects are named without fine-grained
comparison and steering vectors leave no lens-decodable
trace~\cite{shahgir2026vlms,nadaf2026steerable}. Medical lenses track diagnostic formation in a
breast-cancer transformer~\cite{takahashi2026vision}, place multilingual knowledge at intermediate
layers but not at output~\cite{abouzahir2026bridging}, and find early visual attributes regardless
of answer correctness~\cite{nooralahzadeh2026arbitration}.

\subsubsection{Attention, visual tokens and the connector}

Attention to the image is measured as visual-token mass $\sum_{q\in\mathcal V}A^{(\ell,h)}_{pq}$
over heads and layers, or by rollout and its
refinements~\cite{abnar2020quantifying,chefer2021transformer}; none quantifies causal contribution,
as the attention-as-explanation debate
established~\cite{jain2019attention,wiegreffe2019attention,serrano2019is}. General-domain work finds
middle-layer concentration on real objects, staged flow into question tokens, attention sinks,
redundant visual tokens and modality-specific
neurons~\cite{jiang2024devils,zhang2024cross,kang2025you,luo2025sink,chen2024image,wang2025token,huo2024mmneuron,huang2024miner,kang2025large,kim2025interpreting,jiang2025investigating}.
Connectors map correlated embeddings into language-model feature
space~\cite{lee2026more,venhoff2025visual} and hallucinate when
text-over-aligned~\cite{fan2026visual,venhoff2025too,saini2026when}; dermatology fine-tuning
improves the projection but not its domain specificity~\cite{verma2024cross}.

Eleven core records from eleven studies, nine of which take tracing as their primary family (\S\ref{sec:method}), read attention or attribution maps of medical models; \Vtrany{} test them against
perturbation, occlusion or counterfactual evidence, and where tested against a perturbation or counterfactual, the maps often fail, and the result is benchmark- and end-point-dependent. Four attribution methods disagree under patch deletion and insertion across 17 RETFound
models, gradient-weighted rollout winning in 11 of 17~\cite{bors2026benchmarking}, and internal
attributions of six vision--language models miss the counterfactually verified evidence region on
chest radiographs (rollout IoU 2.7 against 39.2 for the best gradient
method)~\cite{xiong2026rethinking}. Attention maps sometimes anti-correlate with occlusion
importance~\cite{sadanandan2026attention}, lose to perturbation and relevance-propagation
explanations~\cite{idaji2026beyond}, enrich for pathway rather than gene
signatures~\cite{srikanthan2026do}, mark pixels outside the body~\cite{hashmi2024towards} and are
diffuse and prompt-insensitive~\cite{injarabian2025interpreting}, while an EchoPrime probe beats
random masking only by a method-dependent margin~\cite{akanova2026interpreting}$^{\dagger}$.

The remaining records localise rather than test: visual-attention ratios find a visually dominant
window in two medical multimodal LLMs (\S\ref{sec:intervention})~\cite{chen2026vihd}, and
class-label probes at successive loci of fourteen models name four failure modes including connector
fidelity loss~\cite{zhu2026lost}. The earliest record audits GLoRIA's text--image attention on Chest
ImaGenome pairs: the map localises above chance (AUROC 69.1) yet barely moves when ``left'' becomes
``right'' and falls only when the reference boxes move, so it is fixed by the
image~\cite{mcinerney2022that}. Text-tower anchors trace negation failure in a chest-radiograph
assistant and in radiology captions~\cite{vu2025improving,abbasi2026layerspecific}; input audits
agree that image use is weak, a text-only baseline nearly matching the best multimodal
model~\cite{lotfinia2026vision}, VQA grounding staying poor~\cite{chen2026grounding} and grounding
failures separating~\cite{liu2026medical}. No other study measures a medical connector
(Fig.~\ref{fig:corpus}b), so these token-level results transfer only where projectors preserve patch
correspondence (Table~\ref{tab:models}).

\subsubsection{Encoder tokens, attention heads and slide-level aggregators}

Three observations bear on the encoders of Table~\ref{tab:models}, none a core claim. The high-norm
artefact tokens of vision transformers~\cite{darcet2023vision} are unmeasured in DINOv2-recipe
models such as RAD-DINO, UNI and Prov-GigaPath, where they would dominate the patch norms and
attention probes read. Self-supervised pathology encoders develop heads attending to morphological
phenotypes~\cite{chen2022self,chen2024generalpurpose,xu2024wholeslide}, qualitative evidence that
morphology was learned. Tile weights of attention-based multiple-instance learning (MIL)
aggregators~\cite{ilse2018attention,lu2020data} are the common computational-pathology explanation,
partly answered by prototype- and concept-based aggregators
(\S\ref{sec:intrinsic})~\cite{song2024morphological,huy2025interactive}. Like decoder attention,
none is confirmed by patching.

\subsubsection{Example- and training-data-based attribution}

Frozen embeddings support content-based retrieval over 1.6 million radiological images~\cite{denner2024leveraging}, and foundation-feature prototypes (\S\ref{sec:intrinsic})
turn nearest-neighbour structure into explanation. Text-based decomposition attributes CLIP
representations to layers, heads and
patches~\cite{gandelsman2023interpreting,balasubramanian2024decomposing}, with no medical instance.
Influence functions~\cite{koh2017influence} trace predictions to training examples, identifying mislabelled
data~\cite{hao2020inaccurate} and re-weighting noisy labels~\cite{braun2022influence}, and scalable
gradient attribution has reached a biomedical language model~\cite{cao2026scalable}; medical vision
use stays limited to curation and leakage auditing (\S\ref{sec:applications}), where image--report
correspondence creates a re-linkage risk~\cite{tayebiarasteh2026crossmodal}.


\subsubsection{Summary and discussion}

Tracing localises information and shows how unevenly it is attended, not that the output uses it: in
language models the layer that causal tracing singles out is not the layer at which editing works
best~\cite{hase2023does}. The medical records that test their maps find them unreliable or
method-dependent, and image use weak. \Vtrany{} of the \Vtrn{} records
\emph{test} their map, by occlusion, perturbation, patching or counterfactual evidence, and the test
corroborates the map in some and refutes it in others; localisation itself remains a hypothesis, and artefact tokens threaten the
patch-level analyses reviewed here.

\subsection{Causal intervention and representation control}
\label{sec:intervention}

The family holds three subtypes that share a name but not an estimand, and their records should not be
read as one total: interventions inside the forward pass, post-hoc transformations of stored embeddings,
and parameter editing or unlearning. Only the first intervenes on an activation inside a fixed model and forward pass; the second changes
what a refitted readout can do with a stored vector; the third changes the model's parameters and
identifies no representation-level mediator unless the study measures one. Each therefore takes a different control
(Table~\ref{tab:families}). Within the forward pass, non-zero effects establish output sensitivity to a
locus, and matched controls direction selectivity. A validated direction strengthens the second, but a
concept-specific contribution needs a mediation or causal-abstraction design, which no reviewed medical
study reports (\S\ref{sec:methods}). Methods include causal mediation and
abstraction, causal tracing, rank-one editing, activation patching and additive
steering~\cite{vig2020causal,geiger2021causal,geiger2023causal,meng2022locating,meng2022mass,zhang2023towards,turner2023steering,rimsky2023steering,zou2023representation,li2023inference,arditi2024refusal}.
Medical claims date from 2025--2026 and remain sparse and decoder-centred (Fig.~\ref{fig:corpus}b),
after early orthogonalisation of chest-radiograph embeddings~\cite{weber2023posthoc} and a
joint-space edit~\cite{sadashivaiah2024explaining}.

\subsubsection{Ablation and knockout}

Ablation removes a component: masking the highest-attention 10\% of visual tokens in the visually
dominant decoder layers of two medical multimodal LLMs detects hallucinations (AUC 71.6--78.9
against 68.9--73.0 for the best baseline), while random masking changes little~\cite{chen2026vihd}.
Ablating object-specific visual tokens likewise reduces object identification in general
models~\cite{neo2024towards}, and knockout in frozen CT encoders shows about ten channels carrying
each radiological finding~\cite{nooralahzadeh2026sparse}. Fine-tuning visually critical heads from
weak segmentation labels is a training method and therefore contextual~\cite{chang2025focus}.

\subsubsection{Activation patching and causal tracing}

Causal tracing corrupts a run, restores clean-run hidden states and measures output recovery; the
indirect effect of locus $\lambda$ in~(\ref{eq:effect}) is
$B(x_{\mathrm{corr}};z_\lambda\!\leftarrow\!z_\lambda(x_{\mathrm{clean}}))-B(x_{\mathrm{corr}})$.
Adaptations to vision--language models localise image-conditioned generation to layers and
tokens~\cite{palit2023towards,basu2024understanding,zhang2024cross,li2025causal}. One medical
instance is core, and it patches a feature rather than tracing a corruption: in MedGemma, patching a
single dictionary feature recovers more of the decision margin than random control
features~\cite{sadanandan2026mechanistically} (F9a). No medical clean--corrupted tracing is core: the
closest scales CLIP text-tower updates by causal contribution~\cite{abbasi2026layerspecific}; text-only anchors place editing
adapters~\cite{xu2024editing}, localise gender~\cite{ahsan2025elucidating} and show knowledge
failing to reach output~\cite{abouzahir2026bridging}. Automated circuit discovery remains
general-domain~\cite{conmy2023automated,syed2023attribution,marks2024sparse,ameisen2025circuit}.

\subsubsection{Steering vectors and representation engineering}

Steering sets $\tilde z_\lambda=z_\lambda+\alpha v$ and measures $\Delta B(\alpha)$
in~(\ref{eq:effect}); $v$ may be a contrastive mean difference, probe normal or dictionary decoder
vector, with an equal-norm random direction or unrelated concept as controls (Table~S5). General
vision encoders support patch-level steering~\cite{joseph2025steering,stevens2025interpretable}, and
CLIP interventions propagate through the projector~\cite{pach2025sparse}.

Evidence is mixed. Per-token steering improves a report-quality composite in three radiology
multimodal LLMs, one transferring zero-shot to a second corpus~\cite{nooralahzadeh2026universal};
Matryoshka dictionary steering is likewise mixed~\cite{bouzid2025insights}. Contrastive activation
addition, with random and reversed controls, shifts pneumonia classification in one of three frozen
chest-radiograph models~\cite{farina2026translating}, moderate modality-gap shifts improve
supervised performance~\cite{restrepo2026cone}, and lesion-mask amplification improves subtle-lesion
detection across five backbones, though on CT it rests on one model without a random
control~\cite{zeng2026easylens}. Text-only anchors divide too: interventions leave a 53-point
probe--action gap~\cite{basu2026interpretability}, hallucination resists
control~\cite{vankadaru2026readable} and race-linked latents are marginally
useful~\cite{ahsan2025can}, yet lightweight steering helps clinical question
answering~\cite{ray2026do}.

General-domain effects are input-dependent, sometimes sign-reversing and
non-identifiable~\cite{tan2024steeringreliability,braun2025understanding,venkatesh2026nonident};
random directions alone raise harmful compliance from 0\% to 1--13\% and move jailbreak success both
ways~\cite{korznikov2025the,li2026analysing}, and function-vector steering dissociates from lens
decodability~\cite{nadaf2026steerable}.

\subsubsection{Post-hoc transformation, editing and unlearning}

Editing changes weights, unlearning removes data influence, erasure removes concepts, and post-hoc
transformations (whitening, harmonisation, projection, transport) change stored embeddings before re-evaluation;
the boundary is the operation, not the paper's genre: an operation on a frozen representation with a downstream
end-point is core (Table~\ref{tab:boundary}), whereas a method that changes the encoder's training or weights
alone is contextual. Editing fails on medical benchmarks: gradient-based editors
generalise but violate locality, memory-based editors preserve locality without generalising
compositionally~\cite{zhu2026evaluating}, and every editor tested fails under single and sequential
editing~\cite{xu2025medmkeb,wen2025multimededit}. Unlearning is
contextual~\cite{guo2026exploring,hardan2025forgetmi,xu2025from}, at best projecting the target
subspace out of language-model and connector weights~\cite{wu2025hierarchy}. Mechanistically guided
low-rank adaptation improves paraphrase consistency~\cite{sadanandan2026mechanistically}, and
editing 1\% of weights injects persistent misinformation into a text-only
anchor~\cite{han2024medical}.

Null-space projection~\cite{ravfogel2020null,ravfogel2022linear}, linear
erasure~\cite{belrose2023leace} and amnesic probing~\cite{elazar2020amnesic} remove linearly
decodable attributes. Projecting chest-radiograph embeddings off the protected-attribute column
space removes attribute effects on pathology logits and shrinks subgroup AUC gaps, at a macro-AUC change of
$-1.4$ to $+1.8$ points across two model--dataset pairs~\cite{weber2023posthoc}. On retinal embeddings it lowers device decodability only from 0.95 to
0.70--0.86 and leaves a three-year device shift in predicted age~\cite{li2026acquisition}.
Residualising demographics leaves clinical-finding prediction unchanged, a test of
dependence~\cite{kim2026encoding}$^{\dagger}$; the joint-space analogue edits along report-derived
directions~\cite{sadashivaiah2024explaining}.

In pathology the target is acquisition: site prediction falls to chance under
ComBat~\cite{murchan2024deep}$^{\dagger}$ and a gradient-reversal
head~\cite{zhang2025hospital}$^{\dagger}$, with biomarker end-points holding; centre-signature
removal raises robustness but not generalisation~\cite{kmen2026robust}; contrastive projection
raises same-specimen agreement~\cite{carloni2025pathology}; an affine correction raises
cross-scanner retrieval~\cite{donle2026featmap}; and whitening degrades downstream performance
(\S\ref{sec:geometry})~\cite{garciahidalgo2026embedding}. Two core joint-space operations do the
same, scoring CPath-CLIP embeddings against text anchors to raise cross-species tumour
transfer~\cite{arora2026semantic} and regularising a text projection head against counterfactual
text~\cite{balogh2026regularization}$^{\dagger}$. Evaluated only downstream, they establish what
they remove, not what the model uses (F3b); adversarial perturbations on frozen embeddings are a
training method~\cite{xu2024apple}, and demographic attributes (\S\ref{sec:decoding}) remain
erasure's target.


\subsubsection{Summary and discussion}

Patching, ablation and steering show outputs to be sensitive to interventions on medical representations,
but concept-directed effects are input- and model-specific: a matched control applies to the eight records that target a specific direction, feature, unit or token,
whatever sub-group heading the study sits under and whether or not the target carries a name (four of the eight target a named concept, F9a), and three report one. The other thirteen harmonise, project, erase or edit
without such a target, share no estimand with them or with each other, and are inventoried as F9b rather than
pooled into a family total; they need locality and retain-set checks (Table~S5), for which
\S\ref{sec:future} proposes a standard. No study in this corpus intervenes on the connector, and
non-identifiability~\cite{korznikov2025the,venkatesh2026nonident} precludes crediting uncontrolled
improvements to concepts.

\subsection{Intrinsic and concept-aligned designs built on foundation models}
\label{sec:intrinsic}

Intrinsic designs constrain part of a representation to be interpretable by construction, a model
property; analysis stays in the preceding families. Eligible designs carry the intrinsic-design tag and are
coded in the supporting family, decoding except for one discovery-coded prototype design.

\textbf{Concept bottlenecks on frozen medical encoders.} Concept bottleneck models (CBMs) predict
human-defined concepts, then labels from concepts alone, $\hat c=g_c(z_\lambda(x))$ and $\hat y=g_y(\hat c)$,
so labels depend on the representation only through inspectable $\hat c$~\cite{koh2020concept}; concept
embedding models relax the bottleneck to a vector per concept to recover accuracy~\cite{zarlenga2022concept}.
Post-hoc and label-free variants attach to pre-trained encoders, language models proposing and vision--language models
scoring concepts~\cite{yuksekgonul2022post,oikarinen2023label,yang2022language}. GPT-4 concepts scored by
a vision--language model give classifiers across eight medical datasets~\cite{yan2023robust}, textbook-
and PubMed-derived concept spaces degrade less than black-box fine-tuning on constructed confound-reversed
splits, a shortcut test rather than a natural external shift~\cite{yang2024textbook}, and stronger
intermediate- than final-layer associations motivate multi-layer bottlenecks~\cite{wang2025mvp}.
Applications reach volumetric, slide and ultrasound
encoders~\cite{rojas2026echopcbm,tushar2026distilling,gao2025evaluating,sun2025label}: frozen-CT ridge
concept heads recover lung-nodule attributes ($R^2\le 0.24$) yet generalise externally better than
embedding probes, and pathology-encoder graph heads feed thirty kidney-cancer concepts to a survival
model. Distribution shift degrades frozen-feature bottlenecks, partly recovered by concept-layer
adaptation~\cite{choi2024adaptive}; VLM-supervised concepts can differ from expert annotations, and
concept accuracy relates weakly to quality~\cite{debole2025if} (\S\ref{sec:eval-what}).

\textbf{Concept-based vision--language models and report-grounded bottlenecks.} Here the concept layer
sits inside a medical VLM, where radiologist-style observation probabilities yield zero-shot
diagnoses~\cite{pellegrini2023xplainer}. Elicited criteria become textual anchors in a VLM's joint
space~\cite{gao2024aligning} or explainable prompts~\cite{bie2024xcoop}, and large VLMs classify from
their own answers to concept questions, a contextual output-level procedure~\cite{patricio2025cbvlm}; a two-step
variant reads dermoscopic concepts off frozen MONET image--text similarity and passes them to a language model,
beating a trained bottleneck on three datasets, with ground-truth concepts adding 5--21 points and BiomedCLIP concepts
at chance~\cite{patricio2025twostep}. Report-derived concepts
in a chest-radiograph joint embedding provide counterfactuals~\cite{sadashivaiah2024explaining}, and
guideline text supports conformant classification in a general CLIP model~\cite{harmanani2026vision}. A radiology foundation model projects chest radiographs into
language-model concept space, trained on more than 0.87M image--report pairs and validated on
physician-annotated cohorts~\cite{han2026clear}.

\textbf{Prototype and part-based models on foundation features.} Case- and part-prototype models explain by
similarity to exemplars~\cite{chen2019thislooks,barnett2021case,kim2021xprotonet}; foundation-model versions
summarise slide patch embeddings into morphological prototypes~\cite{song2024morphological} or
give patch-level prototypes interactive concept meaning~\cite{huy2025interactive}, a self-interpretable
multiple-instance model pairs handcrafted and deep branches~\cite{kapse2023si}, region-additive logistic heads on
frozen BrainIAC and NeuroFM embeddings of anatomically masked volumes match fine-tuning while randomised
embeddings fall to chance~\cite{zhang2026regionfm}, and
attribute-regularised or concept-whitened spaces align axes with clinical
attributes~\cite{cetin2022attri,hou2024concept}, the medical instances of concept whitening~\cite{chen2020concept}. Being feature-space
points and directions, they are evaluated by decoding and discovery.

\textbf{Critiques.} Bottleneck intervention requires labels to depend only on concepts, yet predictors may
ignore relevant regions~\cite{margeloiu2021do} and soft scores leak non-concept information~\cite{havasi2022addressing}. Random concept projections match learned accuracy, so
accuracy cannot certify meaning~\cite{srivastava2024vlg,schmalwasser2026faithfulness};
information-theoretic measures quantify leakage~\cite{parisini2025leakage,makonnen2025measuring}, and
intervenability extends beyond bottlenecks~\cite{marcinkevics2024beyond}. Inconsistent dermoscopy annotations cap
hard-bottleneck accuracy~\cite{napoles2026concept}, and concept accuracy can
conceal a meaningless layer~\cite{han2026loss}. A sparse autoencoder before a linear
bottleneck~\cite{rao2024discover} inherits discovery instability
(\S\ref{sec:discovery}).


\subsubsection{Summary and discussion}

Intrinsic designs make concept layers auditable and intervenable at task accuracy competitive with the same
backbone without the concept layer, where that comparison is reported, but
faithfulness requires valid concepts and controlled leakage, which accuracy cannot show. \Vintrany{} of the
\Vintrn{} intrinsic-design records report a control on the concept layer (Table~\ref{tab:vcov}), an embedding-randomisation
null~\cite{zhang2026regionfm} and the replacement of predicted concepts by ground
truth~\cite{patricio2025twostep}; comparisons of the same frozen encoder with and without an intrinsic
head, and interventions on $\hat c$ with an expert-validated $g_c$ and measured leakage, are rarely
reported together.

\subsection{Comparison across families}
\label{sec:comparison}

\begin{table*}[!tb]
\centering\footnotesize\setlength\tabcolsep{3pt}\renewcommand{\arraystretch}{1.05}
\caption{What each method family licenses, the controls that license more, and where this corpus stops.
Column~2: what the artefact supports on its own, on held-out data above the study-stated reference
(notation of \S\ref{sec:prelim}, Eqs.~(\ref{eq:probe})--(\ref{eq:sae})). Column~3: tests that license a
strictly stronger reading, none a prerequisite (Table~\ref{tab:vcov}). Column~4: the gap that most limits
the family's evidence here, with the record counts of \S\ref{sec:findings}. Geometry is task-linked
geometry only (\S\ref{sec:models}).}
\label{tab:families}
\begin{tabular}{@{}L{2.6cm}L{4.55cm}L{4.15cm}L{4.35cm}@{}}
\toprule
Family & Claim licensed on its own & Stronger reading needs & Where this corpus stops \\
\midrule
\textbf{Structural and geometric} (task-linked) & shape of the space: gap, alignment, effective rank, CKA ($\Phi$ over $\{z_\lambda\}$); no concept & comparison model; random-initialised or shuffled baseline & comparison model \Vcomp{}/\Vcompn{}; reference distribution \Vncal{}/\Vncaln{} \\
\addlinespace[3pt]
\textbf{Supervised concept mapping and decodability} & $c$ recoverable from $z_\lambda$ by decoders in $\mathcal G$ (Eq.~\ref{eq:probe}) & control task or selectivity; shuffled labels; random encoder; fixed $\mathcal G$ across models & permutation null \Vperm{}/\Vpermn{}; random encoder \Vlrn{}/\Vlrnn{}; selectivity not yet measurable (\S\ref{sec:future}) \\
\quad\emph{zero-shot scoring} & $c$ recoverable by a fixed scorer ($\cos(z_v,z_t)$), once validated on held-out labels & prompt-set sensitivity; held-out validation & expert-validated concepts a minority; prompt sensitivity unreported \\
\quad\emph{concept activation vectors} & logit gradient-associated with a fitted direction (TCAV $S_{c,k}$), not recoverability & random-direction null; refit on shuffled labels & refit on shuffled labels: no record \\
\quad\emph{labelled similarity} & representation and concept distances correspond (RSA); no held-out value & permutation null over the labelled sets & permutation null: one record of five \\
\addlinespace[3pt]
\textbf{Feature decomposition and discovery} & a unit's activation tracks the named content ($D$, $a(x)$, Eq.~\ref{eq:sae}) & reconstruction $R^2$; validated names; cross-seed matching & semantic validation \Vsem{}/\Vsemn{}; cross-seed identity: none \\
\addlinespace[3pt]
\textbf{Attribution, localisation and tracing} & $c$ decodable at, or attended from, $(\ell,p,h)$; use unconfirmed & occlusion or patching confirmation; random-component comparison & input-space corroboration \Vinp{}/\Vinpn{}; component perturbation \Vint{}/\Vintn{} \\
\addlinespace[3pt]
\textbf{Representation intervention and control} (within the forward pass) & output sensitive to $z_\lambda$ ($\Delta B$, Eq.~\ref{eq:effect}); with matched controls, selectively to $v_c$ & matched control directions (\S\ref{sec:future}); equal-size random units & matched control \Vspca{}--\Vspc{}/\Vspcn{} targeted records; no dose--response test \\
\quad\emph{post-hoc transformation of stored embeddings} & changes what a refitted readout can do, not the model's computation & sham or alternative transformation; unrelated-attribute control & downstream evaluation only; no sham transformation \\
\quad\emph{parameter editing and unlearning} & behaviour changes as specified; no mediator unless measured & locality and retain-set checks; a measured mediator & no representation-level mediator identified \\
\midrule
\textbf{Intrinsic and concept-aligned designs} (tag) & label depends on $z_\lambda$ only through $\hat c$; prototypes and anchors mean what they are named & leakage test; random concept projections; intervention on $\hat c$; expert labels & concept-layer control \Vintrany{}/\Vintrn{}; leakage rarely measured \\
\bottomrule
\end{tabular}
\end{table*}
\begin{table*}[!tb]
\centering\footnotesize\setlength\tabcolsep{2.5pt}\renewcommand{\arraystretch}{1.05}
\caption{What each family costs and how far this corpus has taken it. \emph{Concept annotation}: whether
the design needs a concept vocabulary, and how much labelling. \emph{Loci}: where the family has been
applied to medical models (parentheses: reached only in general-domain work). \emph{Compute}: beyond a
forward pass. \emph{Readable}: whether a clinician can read the artefact directly. \emph{Stud./rec.}:
core studies whose primary family is this one / records of this family (for decoding,
\Ndecstudnb{}/\Ndecrecnb{} without the \Nbenchdec{} frozen-feature benchmarks). \emph{Groups}: distinct
first authors among those records. \emph{Base.}: records comparing a medical with a general model on the
same data. \emph{Multi.}: more than one model or dataset. The structural row is task-linked geometry only
(\S\ref{sec:models}). Counts come from the coding sheet, the other cells are the authors' reading, and
the \Nintrinsic{} intrinsic-design records are counted in their analytic family (thirteen decoding, one
discovery) and listed again in the tag row.}
\label{tab:famreq}
\begin{tabular}{@{}L{1.86cm}L{1.82cm}L{1.76cm}L{1.33cm}L{1.53cm}C{0.49cm}C{0.89cm}C{0.79cm}C{0.61cm}C{0.61cm}L{2.55cm}@{}}
\toprule
Family & Concept annotation & Loci (medical) & Compute & Readable & 1st yr & Stud./ rec. & Groups & Base. & Multi. & Main failure mode \\
\midrule
Structural and geometric (task-linked) & none & Enc, Txt, Joint, Dec, (Conn) & low & no & 2022 & 21/30 & 27 & 11 & 29 & geometric quantity read as clinical meaning \\
Supervised concept mapping and decodability & required; high (labels) & Enc, Joint & low & partly (scores yes, probes no) & 2022 & 75/80 & 76 & 31 & 72 & decodability read as use \\
Feature decomposition and discovery & none or weak (post-hoc naming) & Enc, Dec, (Conn) & high (activations at scale, training) & yes, once named and validated & 2024 & 20/25 & 25 & 2 & 14 & unstable features and unvalidated names \\
Attribution, localisation and tracing & optional; low--moderate & Enc, Joint, Dec, Path, (Pre) & moderate & partly (maps, often unfaithful) & 2022 & 9/11 & 11 & 2 & 8 & localisation read as causal confirmation \\
Representation intervention and control & required for concept-\hspace{0pt}specific claims; moderate (exemplars) & Enc, Joint, Dec, (Txt, Conn) & mod.--high (repeated generation) & indirectly (behaviour change) & 2023 & 10/21 & 20 & 5 & 18 & uncontrolled, input-\hspace{0pt}dependent effects \\
\midrule
Intrinsic and concept-aligned designs (tag) & required; high & Enc (concept layer), Joint & low--moderate & yes & 2023 & 14 rec. & 13 & 7 & 13 & leakage; concept invalidity \\
\bottomrule
\end{tabular}
\end{table*}

Table~\ref{tab:families} sets out what each family reports and what its artefact licenses;
Table~\ref{tab:famreq} gives what each costs and how far this corpus has taken it, placing the medical
evidence at the observational end (Fig.~\ref{fig:corpus}b). Within this corpus, and under its family-specific eligibility rules, decoding is the only family with a broad base of distinct first authors (76); discovery is two years old and intervention three, and they have 25 and 20 distinct first authors, a proxy for independent groups rather than a measure of the families' maturity in the field. Two asymmetries in Table~\ref{tab:famreq} are easily
missed: intervention gives the claims of greatest practical interest at the highest cost and with the fewest
controlled results, while the cheapest family addresses availability, often mistaken for use; and readability
and validity pull apart, since discovery and intrinsic designs return the artefacts a clinician can read yet
rest on an unvalidated naming or leakage step.

\textbf{Combining families.} Three pipelines recur. \emph{Discover, then intervene} uses residual-stream
dictionary directions for steering or
patching~\cite{bouzid2025insights,nooralahzadeh2026universal,sadanandan2026mechanistically}, carrying the
caveats of \S\ref{sec:discovery} into the intervention. \emph{Trace, then intervene} exploits an
observational localisation through a training-time change, so far only in a general CLIP text tower on
radiology captions~\cite{abbasi2026layerspecific}. \emph{Decode, then erase, residualise or ablate} probes
for availability, then tests whether a downstream prediction depends on the
concept~\cite{kim2026encoding,nooralahzadeh2026sparse,weber2023posthoc}, answering the shortcut question of
\S\ref{sec:applications} most directly. Of its three core instances the two that test a downstream
dependence point in opposite directions, since
residualising a demographic attribute leaves a refitted finding probe unchanged~\cite{kim2026encoding}$^{\dagger}$
whereas knocking out the roughly ten channels that decode a finding from a CT encoder collapses its
score~\cite{nooralahzadeh2026sparse}; the third orthogonalises chest-radiograph embeddings and reports a
downstream change within noise~\cite{weber2023posthoc}.

Reading one family's result as an answer to another family's question is the commonest inferential error in
this literature, and the evidence-type tag makes it visible.


